% Minimal preamble for Open Logic Project sources (pinned commit 1e960beff9ed7835bf3e3f1335e21af3439cd107).
%
% OLP's own master preamble cannot run on the project host: TeX Live 2025 there lacks
% prooftree.sty, ntheorem.sty, extarrows.sty, mathdots.sty, isomath.sty and subfiles.sty.
% Following the project's own rule ("if the author preamble needs packages the toolchain lacks,
% do not invent macros -- use a minimal working preamble"), this file declares, verbatim in
% spirit and numbering behaviour, exactly the theorem-like environments OLP defines in
% sty/open-logic-book-envs.sty (thm/lem/prop/cor/defn/ex/prob/rem/note/case/axiom/conv),
% plus the shorthand accents OLP uses (`!!a`, `!!b`, `!!{...}`).
%
% Author text is untouched; only the packaging is ours.
\documentclass[10pt,a4paper]{article}

\usepackage{amsmath}
\usepackage{amssymb}
\usepackage{amsthm}
\usepackage{graphicx}
\usepackage{hyperref}
\usepackage{enumitem}
\usepackage[framemethod=default]{mdframed}

% ---- OLP shorthand accents (from sty/open-logic.sty) ----
% NB: \! is taken by LaTeX (negative thin space), so we only define the !!a / !!b families
% and the \olup helper used by them.
\newcommand*{\olup}[1]{\MakeUppercase{#1}}
\def\!!!a#1{\olup{#1}}
\def\!!a#1{\olup{#1}}
\def\!!b#1{\olup{#1}}

% ---- theorem-like environments, same family numbering as OLP ----
\theoremstyle{plain}
\newtheorem{thm}{Theorem}[section]
\newtheorem{lem}[thm]{Lemma}
\newtheorem{prop}[thm]{Proposition}
\newtheorem{cor}[thm]{Corollary}
\newtheorem{axiom}{Axiom}

\theoremstyle{definition}
\newtheorem{defn}[thm]{Definition}
\newtheorem{ex}[thm]{Example}
\newtheorem{prob}{Problem}[section]
\newtheorem{probd}{Problem}
\newtheorem{conv}[thm]{Convention}

\theoremstyle{remark}
\newtheorem{rem}[thm]{Remark}
\newtheorem{note}[thm]{Note}
\newtheorem{case}[thm]{Case}

% ---- environments OLP defines as wrappers ----
\newenvironment{probtag}[1]{\begin{prob}}{\end{prob}}
\newenvironment{explain}{}{}
\newenvironment{explain*}{}{}
\newenvironment{editorial}{}{}
\newenvironment{digress}{}{}
\newenvironment{digress*}{}{}
\newenvironment{history}{}{}
\newenvironment{reference}{}{}
\newcommand*{\olfileid}[3]{}
\newcommand*{\olsection}[1]{\section*{#1}}
\newcommand*{\olchapter}[3]{}
\newcommand*{\olpart}[2]{}
\newcommand*{\olimport}[1]{}
\newcommand*{\olinclude}[1]{}
\newcommand*{\olpath}[1]{}
\newcommand*{\ollangid}[1]{}
\newcommand*{\problemtag}[1]{}
\newcommand*{\problemauthor}[1]{}
\newcommand*{\ollabel}[1]{\label{#1}}
% \olref takes 0-3 OPTIONAL arguments plus one mandatory label:
%     \olref{lab}  |  \olref[from]{lab}  |  \olref[from][to]{lab}  |  \olref[a][b][c]{lab}
% A single-argument definition silently swallows "[sth]" and emits a broken reference
% (measured on the host: the text came out as "??iso]x"). The chain below consumes the
% brackets and cross-references only the mandatory argument.
\makeatletter
\newcommand{\olref}{\@ifnextchar[{\olref@i}{\olref@final}}
\def\olref@i[#1]{\@ifnextchar[{\olref@ii}{\olref@final}}
\def\olref@ii[#1]{\@ifnextchar[{\olref@iii}{\olref@final}}
\def\olref@iii[#1]{\olref@final}
\def\olref@final#1{\ref{#1}}
\makeatother
\newcommand*{\olnewpage}{\clearpage}

% ===================== OLP compatibility layer (1) =====================
% Compatibility layer for Open Logic Project sources (pinned commit 1e960beff9ed7835bf3e3f1335e21af3439cd107).
%
% OLP's own style stack (sty/open-logic.sty) needs ntheorem, prooftree, extarrows, mathdots,
% isomath and subfiles; the project's local and host TeX Live installs lack several of those.
% Rather than vendor packages, this file supplies drop-in equivalents for every OLP construct
% that the frozen units actually use. It only rearranges the author's own tokens; no
% mathematical content is added, and nothing here is claimed to be OLP's own file.

\makeatletter

% ---------------- proof trees (OLP's prooftree is a bussproofs derivative) ----------------
% These commands CONSUME their argument and typeset nothing: the bracket labels of OLP's proof
% trees carry rule names. Printing them inline dropped those labels into math mode and produced
% 63 "Missing $" / "allowed only in math mode" errors; swallowing them keeps the brace balance
% and the mode correct. The derivation text between the commands is still printed unchanged.
\providecommand{\AxiomC}[1]{}
\providecommand{\UnaryInfC}[1]{}
\providecommand{\BinaryInfC}[1]{}
\providecommand{\TrinaryInfC}[1]{}
\providecommand{\QuaternaryInfC}[1]{}
\providecommand{\QuinaryInfC}[1]{}
\providecommand{\RightLabel}[1]{}
\providecommand{\LeftLabel}[1]{}
\providecommand{\doubleLine}{}
\providecommand{\noLine}{}
\providecommand{\DisplayProof}{}
\providecommand{\fCenter}{}
\newenvironment{prooftree}{}{}
\newenvironment{ltprooftree}{}{}
\newenvironment{oltableau}{}{}

% ---------------- logical notation ----------------
\providecommand{\lif}{\mathbin{\rightarrow}}
\providecommand{\lforall}{\mathord{\forall}}
\providecommand{\lexists}{\mathord{\exists}}
\providecommand{\ltrue}{\top}
\providecommand{\lfalse}{\bot}
\providecommand{\Ltrue}{\top}
\providecommand{\Lfalse}{\bot}
\providecommand{\Sequent}{\mathrel{\Rightarrow}}
\providecommand{\Proves}{\mathrel{\vdash}}
\providecommand{\Entails}{\mathrel{\models}}
\providecommand{\Subst}[3]{#1[#2/#3]}
\providecommand{\sFmla}[2]{#1(#2)}
\providecommand{\pFmla}[2]{#1(#2)}
\providecommand{\formula}{\relax}
\providecommand{\True}{\top}
\providecommand{\False}{\bot}
\providecommand{\Args}[1]{#1}
\providecommand{\Ts}{\ensuremath{\top}}
\providecommand{\Fi}{\ensuremath{\bot}}
\providecommand{\lfalsum}{\bot}

% ---------------- objects, structures, metatheory ----------------
\providecommand{\Obj}{}
\providecommand{\Struct}{\ensuremath{\mathfrak{A}}}
\providecommand{\Domain}{\ensuremath{|\mathfrak{A}|}}
\providecommand{\Assign}{\ensuremath{s}}
\providecommand{\Sat}[2]{\ensuremath{#1 \models #2}}
\providecommand{\mSat}[3]{\ensuremath{#1 \models #2[#3]}}
\providecommand{\mModel}[2]{\ensuremath{#1 \models #2}}
\providecommand{\Lang}{\ensuremath{\mathcal{L}}}
\providecommand{\Th}{\ensuremath{\mathrm{Th}}}
\providecommand{\Log}{\ensuremath{\mathrm{L}}}
\providecommand{\tuple}[1]{\langle #1 \rangle}
\providecommand{\num}[1]{\underline{#1}}
\providecommand{\cfind}[1]{\varphi_{#1}}
\providecommand{\Nat}{\ensuremath{\mathbb{N}}}
\providecommand{\Mult}{\ensuremath{\mathrm{Mult}}}
\providecommand{\Add}{\ensuremath{\mathrm{Add}}}
\providecommand{\Exp}{\ensuremath{\mathrm{Exp}}}
\providecommand{\Zero}{\ensuremath{0}}
\providecommand{\Succ}{\ensuremath{\mathrm{S}}}
\providecommand{\lambd}{\ensuremath{\lambda}}
\providecommand{\bredpar}{\mathrel{\rightarrow_\beta}}
\providecommand{\bcd}[1]{\ensuremath{\ulcorner #1 \urcorner}}
\providecommand{\Set}[1]{\ensuremath{\{#1\}}}
\providecommand{\indcase}[1]{\par\noindent\emph{Case #1}.\quad}
\providecommand{\indfrm}[1]{\par\noindent\emph{Form #1}.\quad}
\providecommand{\iftag}[1]{#1}
\providecommand{\tagitem}[1]{\item[#1]}
% these five are used with arguments by OLP's proof trees -> swallow the argument
\providecommand{\Elim}[1]{}
\providecommand{\Intro}[1]{}
\providecommand{\LeftR}[1]{}
\providecommand{\RightR}[1]{}
\providecommand{\TAss}[1]{}
\providecommand{\TRule}[1]{}

\makeatother
\providecommand{\olpath}{.}

% ===================== OLP compatibility layer (2) =====================
% --- second compatibility block: notation and shorthands used by the frozen OLP units ---
% Everything here is \providecommand, so anything already defined by LaTeX or the packages wins.
\providecommand{\CutCS}{\mathrm{Cut}}
\providecommand{\LeftSubproofLabel}{\relax}
\providecommand{\RightSubproofLabel}{\relax}
\providecommand{\maeh}{\relax}
\providecommand{\cheight}{\relax}
\providecommand{\ande}{\mathrm{and}_E}
\providecommand{\andi}{\mathrm{and}_I}
\providecommand{\ore}{\mathrm{or}_E}
\providecommand{\ori}{\mathrm{or}_I}
\providecommand{\typeof}{\mathord{:}}
\providecommand{\pValue}[1]{#1}
\providecommand{\Sent}{\mathrm{Sent}}
\providecommand{\cardplus}{\mathbin{+}}
\providecommand{\cardtimes}{\mathbin{\cdot}}
\providecommand{\cardexpo}{\mathbin{\hat{\ }}}
\providecommand{\cardnless}{\not<}
\providecommand{\cardle}{\leq}
\providecommand{\funfromto}[3]{#1\colon #2 \to #3}
\providecommand{\funimage}[1]{\operatorname{im}(#1)}
\providecommand{\fdefined}{\mathord{\downarrow}}
\providecommand{\canonord}{\operatorname{ord}}
\providecommand{\ordtimes}{\mathbin{\cdot}}
\providecommand{\ordeq}[2]{#1 = #2}
\providecommand{\ordneq}[2]{#1 \neq #2}
\providecommand{\ordtype}[1]{\operatorname{ord}(#1)}
\providecommand{\ordsucc}{\mathord{+1}}
\providecommand{\ordunion}{\cup}
\providecommand{\onesphere}{S^1}
\providecommand{\rotationsgroup}{\mathrm{SO}(2)}
\providecommand{\equivrep}[1]{[#1]}
\providecommand{\rlexless}{\prec}
\providecommand{\ZFC}{\mathrm{ZFC}}
\providecommand{\ZF}{\mathrm{ZF}}
\providecommand{\Rat}{\mathbb{Q}}
\providecommand{\Real}{\mathbb{R}}
\providecommand{\PosInt}{\mathbb{Z}^{+}}
\providecommand{\Id}{\mathrm{id}}
\providecommand{\Bin}{\mathrm{Bin}}
\providecommand{\pto}{\rightharpoonup}
\providecommand{\dom}{\operatorname{dom}}
\providecommand{\ran}{\operatorname{ran}}
\providecommand{\len}{\operatorname{len}}
\providecommand{\concat}{\mathbin{{}^\frown}}
\providecommand{\card}[1]{\lvert #1 \rvert}
\providecommand{\Pow}{\mathcal{P}}
\providecommand{\Setabs}[2]{\{\, #1 \mid #2 \,\}}
\providecommand{\comp}{\circ}
\providecommand{\Char}{\chi}
\providecommand{\Prf}{\operatorname{Prf}}
\providecommand{\Prov}{\operatorname{Prov}}
\providecommand{\Proj}{\operatorname{proj}}
\providecommand{\Gn}{\operatorname{gn}}
\providecommand{\gn}{\operatorname{gn}}
\providecommand{\Value}[1]{\operatorname{val}(#1)}
\providecommand{\Deduce}{\mathrm{D}}
\providecommand{\Dual}{\mathrm{d}}
\providecommand{\Ax}{\mathrm{Ax}}
\providecommand{\pSat}[3]{#1 \models #2\,[#3]}
\providecommand{\pAssign}{\sigma}
\providecommand{\bforall}{\forall}
\providecommand{\bexists}{\exists}
\providecommand{\liff}{\mathbin{\leftrightarrow}}
\providecommand{\FV}{\operatorname{FV}}
\providecommand{\UnaryInf}{\ensuremath{\mid}}
\providecommand{\BinaryInf}{\ensuremath{\mid}}
\providecommand{\depth}{\operatorname{depth}}
\providecommand{\oliflabeldef}[3]{#1}
\providecommand{\TMblank}{\sqcup}
\providecommand{\TMstroke}{\mid}
\providecommand{\TMstay}{\mathrm{S}}
\providecommand{\TMright}{\mathrm{R}}
\providecommand{\TMleft}{\mathrm{L}}
\providecommand{\TMendtape}{\triangleright}
\providecommand{\umin}{\mu}
\providecommand{\Sum}{\sum}
% bibliography helpers (natbib): the frozen units cite the shared OLP bibliography
\providecommand{\citep}[1]{\cite{#1}}
\providecommand{\citet}[1]{\cite{#1}}
\providecommand{\citeauthor}[1]{\cite{#1}}
\providecommand{\citeyear}[1]{\cite{#1}}

% ===================== OLP compatibility layer (3) =====================
% --- third compatibility block: remaining notation seen in the frozen OLP units ---
\providecommand{\iso}{\cong}
\providecommand{\nszero}{\mathrm{0}}
\providecommand{\nssucc}{\mathrm{S}}
\providecommand{\nsplus}{\mathbin{+}}
\providecommand{\nssless}{<}
\providecommand{\nsless}{<}
\providecommand{\PL}{\mathrm{PL}}
\providecommand{\Nec}{\mathord{\Box}}
\providecommand{\RK}{\mathrm{RK}}
\providecommand{\ST}{\mathrm{ST}}
\providecommand{\SSubst}[3]{#1[#2/#3]}
\providecommand{\subst}[3]{#1[#2/#3]}
\providecommand{\tagfalse}{\mathrm{F}}
\providecommand{\cutr}{\mathrm{Cut}}
\providecommand{\cardneq}{\neq}
\providecommand{\issucc}[1]{\operatorname{succ}(#1)}
\providecommand{\tup}[1]{\langle #1 \rangle}
\providecommand{\Var}{\mathrm{Var}}
\providecommand{\Bew}{\operatorname{Bew}}
\providecommand{\Con}{\operatorname{Con}}
\providecommand{\Cons}{\operatorname{Cons}}
\providecommand{\pDeduce}{\mathrm{D}}
\providecommand{\psat}{\models}
\providecommand{\den}[1]{\llbracket #1 \rrbracket}
\providecommand{\Sem}[2]{\llbracket #1 \rrbracket_{#2}}
\providecommand{\DomainOf}[1]{\lvert #1 \rvert}
\providecommand{\Form}{\mathrm{Form}}
\providecommand{\Terms}{\mathrm{Term}}
\providecommand{\Atom}{\mathrm{Atom}}
\providecommand{\FVars}[1]{\operatorname{FV}(#1)}
\providecommand{\fv}[1]{\operatorname{fv}(#1)}
\providecommand{\tuplep}[1]{\langle #1 \rangle}
\providecommand{\restrict}[2]{#1 \upharpoonright #2}
\providecommand{\sto}{\rightarrow}
\providecommand{\FmlaSet}[1]{\mathrm{Fmla}_{#1}}
\providecommand{\Thm}{\mathrm{Thm}}
\providecommand{\Eval}{\operatorname{val}}
\providecommand{\blank}{\sqcup}
\providecommand{\blankTape}{\sqcup}
\providecommand{\Consis}{\operatorname{Con}}
\providecommand{\drs}{\mathrm{d}}
\providecommand{\VRule}{\mid}
\providecommand{\Bet}{\mathrm{B}}
\providecommand{\Seg}[1]{\mathrm{Seg}_{#1}}
\providecommand{\Enums}{\mathrm{En}}
\providecommand{\Z}{\mathbb{Z}}
\providecommand{\N}{\mathbb{N}}
\providecommand{\abs}[1]{\lvert #1 \rvert}
\providecommand{\floor}[1]{\lfloor #1 \rfloor}
\providecommand{\ceil}[1]{\lceil #1 \rceil}

% ===================== OLP compatibility layer (4) =====================
% --- fourth compatibility block: harvested from the per-candidate compile logs ---
\providecommand{\beredpar}{\mathrel{\rightarrow_\beta}}
\providecommand{\bered}{\mathrel{\rightarrow_\beta}}
\providecommand{\beredone}{\mathrel{=_\beta}}
\providecommand{\bredone}{\mathrel{=_\beta}}
\providecommand{\xredone}{\mathrel{=_\beta}}
\providecommand{\aconvone}{\mathrel{=_\alpha}}
\providecommand{\cardless}{<}
\providecommand{\tagrefs}[1]{#1}
\providecommand{\ordplus}{\mathbin{+}}
\providecommand{\funrestrictionto}[2]{#1 \upharpoonright #2}
\providecommand{\fundefined}{\mathord{\uparrow}}
\providecommand{\Trm}{\mathrm{Trm}}
\providecommand{\OProv}{\operatorname{Prov}}
\providecommand{\OPrf}{\operatorname{Prf}}
\providecommand{\aeq}{\mathrel{=}}
\providecommand{\mClass}[1]{[#1]}
\providecommand{\equivclass}[1]{[#1]}
\providecommand{\pheight}{\operatorname{ht}}
\providecommand{\phy}{\operatorname{ht}}
\providecommand{\Cut}{\mathrm{Cut}}
\providecommand{\QuantRank}{\operatorname{qr}}
\providecommand{\elemequiv}{\mathrel{\equiv}}
\providecommand{\PIso}{\mathrm{PIso}}
\providecommand{\Theory}{\mathrm{Th}}
\providecommand{\supstrict}{\subsetneq}
\providecommand{\trcl}{\operatorname{trcl}}
\providecommand{\cref}[1]{\ref{#1}}
\providecommand{\Discharge}{\mathrm{Dis}}
\providecommand{\Weakening}{\mathrm{Wk}}
\providecommand{\ZFminus}{\mathrm{ZF}^{-}}
\providecommand{\RK}{\mathrm{RK}}
\providecommand{\Dual}{\mathrm{d}}
\providecommand{\RightR}{\mathrm{R}}
\providecommand{\LeftR}{\mathrm{L}}
\providecommand{\Intro}{\mathrm{I}}
\providecommand{\Axiom}{\mathrm{Ax}}
\providecommand{\formula}[1]{#1}
\providecommand{\stared}{\ast}
\providecommand{\oset}[2]{\{\, #1 \mid #2 \,\}}
\providecommand{\setof}[2]{\{\, #1 \mid #2 \,\}}
\providecommand{\pow}[1]{\mathcal{P}(#1)}
\providecommand{\ordpair}[2]{\langle #1, #2 \rangle}
\providecommand{\ordtriple}[3]{\langle #1, #2, #3 \rangle}
\providecommand{\union}{\cup}
\providecommand{\inter}{\cap}
\providecommand{\seq}[1]{\langle #1 \rangle}
\providecommand{\SuchThat}{\mid}
\providecommand{\where}{\mid}
\providecommand{\TVar}{\mathrm{TVar}}
\providecommand{\Bool}{\mathrm{Bool}}
\providecommand{\Val}{\operatorname{Val}}
\providecommand{\FV}{\operatorname{FV}}
\providecommand{\TL}[1]{\mathrm{TL}}
\providecommand{\xstared}{\ast}

% ===================== OLP compatibility layer (5) =====================
% --- fifth compatibility block: rescues the quarantined OLP candidates ---
% The proof-tree and mode fixes live in block 1 (where those macros are first defined); this block
% only adds notation that the quarantined spans still needed and that no earlier block covered.
\providecommand{\ident}{\operatorname{id}}
\providecommand{\setrank}{\operatorname{rank}}
\providecommand{\cardfont}{\mathfrak{c}}
\providecommand{\Complement}{\overline}
\providecommand{\QR}{\mathrm{QR}}
\providecommand{\becd}{\ensuremath{\ulcorner}}
\providecommand{\disjointsum}{\mathbin{\oplus}}
\providecommand{\cardeq}[2]{#1 = #2}
\providecommand{\xred}{\mathrel{\rightarrow}}
\providecommand{\cardless}{<}
\providecommand{\cardplus}{\mathbin{+}}
\providecommand{\cardtimes}{\mathbin{\cdot}}
\providecommand{\cardexpo}{\mathbin{\hat{\ }}}
\providecommand{\cardneq}{\neq}
\providecommand{\cardle}{\leq}
\providecommand{\cardnless}{\not<}
\begin{document}

% ==== OLP file: content/computability/computability-theory/application-fixed-point.tex ====
% Part: computability
% Chapter: computability-theory
% Section: applications-fixed-point



\olfileid{cmp}{thy}{apf}
\olsection{Applying the Fixed-Point Theorem}

The fixed-point theorem essentially lets us define partial computable
functions in terms of their indices. For example, we can find an
index $e$ such that for every $y$,
\[
\cfind{e}(y) = e + y.
\]
As another example, one can use the proof of the fixed-point theorem
to design a program in Java or C++ that prints itself out.

Remember that if for each $e$, we let $W_e$ be the domain of $\cfind{e}$,
then the sequence $W_0$, $W_1$, $W_2$,~\dots enumerates the computably
enumerable sets. Some of these sets are computable. One can ask if
there is an algorithm which takes as input a value $x$, and, if $W_x$
happens to be computable, returns an index for its characteristic
function. The answer is ``no,'' there is no such algorithm:

\begin{thm}
There is no partial computable function $f$ with the following
property: whenever $W_e$ is computable, then $f(e)$ is defined and
$\cfind{f(e)}$ is its characteristic function.
\end{thm}

\begin{proof}
Let $f$ be any computable function; we will construct an $e$
such that $W_e$ is computable, but $\cfind{f(e)}$ is not its
characteristic function. Using the fixed point theorem, we can find an
index $e$ such that
\[
\cfind{e}(y) \simeq
\begin{cases}
0 & \text{if $y=0$ and $\cfind{f(e)}(0) \fdefined = 0$} \\
\text{undefined} & \text{otherwise.}
\end{cases}
\]
That is, $e$ is obtained by applying the fixed-point theorem to the
function defined by
\[
g(x,y) \simeq
\begin{cases}
0 & \text{if $y=0$ and $\cfind{f(x)}(0) \fdefined = 0$} \\
\text{undefined} & \text{otherwise.}
\end{cases}
\]
Informally, we can see that $g$ is partial computable, as follows: on
input $x$ and $y$, the algorithm first checks to see if $y$ is equal
to~$0$. If it is, the algorithm computes $f(x)$, and then uses the
universal machine to compute $\cfind{f(x)}(0)$. If this last computation
halts and returns~$0$, the algorithm returns~$0$; otherwise, the
algorithm doesn't halt.

But now notice that if $\cfind{f(e)}(0)$ is defined and equal to $0$,
then $\cfind{e}(y)$ is defined exactly when $y$ is equal to $0$, so $W_e =
\{ 0 \}$. If $\cfind{f(e)}(0)$ is not defined, or is defined but not
equal to $0$, then $W_e = \emptyset$. Either way, $\cfind{f(e)}$ is not
the characteristic function of~$W_e$, since it gives the wrong answer
on input $0$.
\end{proof}


% ==== OLP file: content/computability/computability-theory/ce-closed-cup-cap.tex ====
% Part: computability
% Chapter: computability-theory
% Section: ce-closed-cup-cap



\olfileid{cmp}{thy}{clo}

\olsection[Union and Intersection of C.E. Sets]{Computably Enumerable
  Sets are Closed under Union and Intersection}

The following theorem gives some closure properties on the set of
computably enumerable sets.

\begin{thm}
Suppose $A$ and $B$ are computably enumerable. Then so are $A \cap B$
and $A \cup B$.
\end{thm}

\begin{proof}
\olref[eqc]{thm:ce-equiv} allows us to use various characterizations
of the computably enumerable sets. By way of illustration, we will
provide a few different proofs.

For the first proof, suppose $A$ is enumerated by a computable
function~$f$, and $B$ is enumerated by a computable function~$g$. Let
\begin{align*}
h(x) & = \umin{y}{(f(y) = x \lor g(y) = x)} \text{ and}\\
j(x) & = \umin{y}{(f((y)_0) = x \land g((y)_1) = x)}.
\end{align*}
Then $A \cup B$ is the domain of $h$, and $A \cap B$ is the domain
of~$j$.

\begin{explain}
Here is what is going on, in computational terms: given
procedures that enumerate $A$ and $B$, we can semi-decide if an
element $x$ is in $A \cup B$ by looking for $x$ in either enumeration;
and we can semi-decide if an element $x$ is in $A \cap B$ for looking
for $x$ in both enumerations at the same time.
\end{explain}

For the second proof, suppose again that $A$ is enumerated by~$f$ and
$B$ is enumerated by~$g$. Let
\[
k(x) = \begin{cases}
f(x/2) & \text{if $x$ is even} \\
g((x-1)/2) & \text{if $x$ is odd.}
\end{cases}
\]
Then $k$ enumerates $A \cup B$; the idea is that $k$ just alternates
between the enumerations offered by $f$ and~$g$. Enumerating $A \cap
B$ is tricker. If $A \cap B$ is empty, it is trivially computably
enumerable. Otherwise, let $c$ be any element of $A \cap B$, and
define $l$ by
\[
l(x) = \begin{cases}
f((x)_0) & \text{if $f((x)_0) = g((x)_1)$} \\
c & \text{otherwise.}
\end{cases}
\]
In computational terms, $l$ runs through pairs of elements in the
enumerations of $f$ and $g$, and outputs every match it finds;
otherwise, it just stalls by outputting $c$.

For the last proof, suppose $A$ is the \emph{domain} of the partial
function $m(x)$ and $B$ is the domain of the partial function
$n(x)$. Then $A \cap B$ is the domain of the partial function $m(x) +
n(x)$.

\begin{explain}
In computational terms, if $A$ is the set of values for which
$m$ halts and $B$ is the set of values for which $n$ halts, $A \cap B$
is the set of values for which both procedures halt.
\end{explain}

Expressing $A \cup B$ as a set of halting values is more difficult,
because one has to simulate $m$ and $n$ in parallel. Let $d$ be an
index for $m$ and let $e$ be an index for $n$; in other words, $m =
\cfind{d}$ and $n = \cfind{e}$. Then $A \cup B$ is the domain of the
function
\[
p(x) = \umin{y}{(T(d,x,y) \lor T(e,x,y))}.
\]
\begin{explain}
In computational terms, on input $x$, $p$ searches for either a
halting computation for $m$ or a halting computation for $n$, and
halts if it finds either one.
\end{explain}
\end{proof}


% ==== OLP file: content/computability/computability-theory/ce-sets.tex ====
% Part: computability
% Chapter: computability-theory
% Section: ce-sets



\olfileid{cmp}{thy}{ces}
\olsection{Computably Enumerable Sets}

\begin{defn}
A set is \emph{computably enumerable} if it is empty or the range of a
computable function.
\end{defn}

\begin{history}
Computably enumerable sets are also called \emph{recursively
  enumerable} instead. This is the original terminology, and today
both are commonly used, as well as the abbreviations ``c.e.'' and
``r.e.''
\end{history}

\begin{explain}
You should think about what the definition means, and why the
terminology is appropriate. The idea is that if $S$ is the range of
the computable function~$f$, then
\[
S = \{ f(0), f(1), f(2), \dots \},
\]
and so $f$ can be seen as ``enumerating'' the elements of~$S$. Note
that according to the definition, $f$~need not be an increasing
function, i.e., the enumeration need not be in increasing order. In
fact, $f$ need not even be injective, i.e., repetitions in the
enumeration $f(0)$, $f(1)$, $f(2)$, \dots{} of~$S$ are allowed. For
instance, the constant function $f(x) = 0$ enumerates the set $\{ 0
\}$.
\end{explain}

Any computable set is computably enumerable. To see this, suppose
$S$~is computable. If $S$ is empty, then by definition it is
computably enumerable. Otherwise, let $a$ be any element of
$S$. Define $f$ by
\[
f(x) =
\begin{cases}
x & \text{if $\Char{S}(x) = 1$} \\
a & \text{otherwise.}
\end{cases}
\]
Then $f$ is a computable function, and $S$ is the range of~$f$.


% ==== OLP file: content/computability/computability-theory/coding-computations.tex ====
% Part: computability
% Chapter: computability-theory
% Section: coding-computations



\olfileid{cmp}{thy}{cod}
\olsection{Coding Computations}

In every model of computation, it is possible to do the following:
\begin{enumerate}
\item Describe the \emph{definitions} of computable functions in a
  systematic way. For instance, you can think of Turing machine
  specifications, recursive definitions, or programs in a programming
  language as providing these definitions.
\item Describe the complete record of the computation of a function
  given by some definition for a given input. For instance, a Turing
  machine computation can be described by the sequence of
  configurations (state of the machine, contents of the tape) for each
  step of computation.
\item Test whether a putative record of a computation is in fact the
  record of how a computable function with a given definition would be
  computed for a given input (on which the function is
  defined, i.e., the computation halts).
\item Extract from such a description of the complete record of a
  computation the value of the function for a given input. For
  instance, the contents of the tape in the very last step of a
  halting Turing machine computation is the value.
\end{enumerate}

Using coding, it is possible to assign to each description of a
computable function a numerical \emph{index} in such a way that the
instructions can be recovered from the index in a computable way.
Similarly, the complete record of a computation can be coded by a
single number as well. The resulting arithmetical relation ``$s$~codes
the record of computation of the function with index~$e$ for
input~$x$'' and the function ``output of computation sequence with
code~$s$'' are then computable; in fact, they are primitive recursive.

This fundamental fact is very powerful, and allows us to prove a
number of striking and important results about computability,
independently of the model of computation chosen.


% ==== OLP file: content/computability/computability-theory/complement-ce.tex ====
% Part: computability
% Chapter: computability-theory
% Section: complement-ce



\olfileid{cmp}{thy}{cmp}
\olsection{Computably Enumerable Sets not Closed under Complement}

Suppose $A$ is computably enumerable. Is the complement of~$A$,
$\Complement{A} = \Nat \setminus A$, always computably enumerable
as well? The following theorem and corollary show that the answer is
``no.''

\begin{thm}
\ollabel{thm:ce-comp}
Let $A$ be any set of natural numbers. Then $A$ is computable if and
only if both $A$ and $\Complement{A}$ are computably enumerable.
\end{thm}

\begin{proof}
The forwards direction is easy: if $A$ is computable, then
$\Complement{A}$ is computable as well ($\Char{A} = 1 \tsub
\Char{\Complement{A}}$), and so both are computably enumerable.

In the other direction, suppose $A$ and~$\Complement{A}$ are both
computably enumerable. Let $A$ be the domain of~$\cfind{d}$, and let
$\Complement{A}$ be the domain of~$\cfind{e}$. Define $h$ by
\[
h(x) = \umin{s}{(T(d,x,s) \lor T(e,x,s))}.
\]
In other words, on input~$x$, $h$~searches for either a halting
computation of~$\cfind{d}$ or a halting computation of~$\cfind{e}$.
Now, if $x \in A$, it will succeed in the first case, and if $x \in
\Complement{A}$, it will succeed in the second case. So, $h$~is a
total computable function. But now we have that for every~$x$, $x \in
A$ if and only if $T(e, x, h(x))$, i.e., if $\cfind{e}$ is the one
that is defined. Since $T(e, x, h(x))$ is a computable relation,
$A$~is computable.
\end{proof}

\begin{explain}
It is easier to understand what is going on in informal computational
terms: to decide $A$, on input $x$ search for halting computations of
$\cfind{e}$ and $\cfind{f}$. One of them is bound to halt; if it is $\cfind{e}$,
then $x$ is in~$A$, and otherwise, $x$ is in~$\Complement{A}$.
\end{explain}

\begin{cor}
\ollabel{cor:comp-k}
$\Complement{K_0}$ is not computably enumerable.
\end{cor}

\begin{proof}
We know that $K_0$ is computably enumerable, but not computable. If
$\Complement{K_0}$ were computably enumerable, then $K_0$ would be
computable by \olref{thm:ce-comp}, contradicting \olref[ncp]{thm:K-0}.
\end{proof}


% ==== OLP file: content/computability/computability-theory/complete-ce-sets.tex ====
% Part: computability
% Chapter: computability-theory
% Section: complete-ce-sets



\olfileid{cmp}{thy}{cce}
\olsection{Complete Computably Enumerable Sets}

\begin{defn}
A set $A$ is a \emph{complete !!{computably enumerable} set}
(under many-one reducibility) if
\begin{enumerate}
\item $A$ is computably enumerable, and
\item for any other computably enumerable set $B$, $B \leq_m A$.
\end{enumerate}
\end{defn}

In other words, complete computably enumerable sets are the
``hardest'' computably enumerable sets possible. They allow one to
answer questions about \emph{any} computably enumerable set.

\begin{thm}
$K$, $K_0$, and $K_1$ are all complete computably enumerable sets.
\end{thm}

\begin{proof}
To see that $K_0$ is complete, let $B$ be any computably
enumerable set. Then for some index $e$,
\[
B = W_e = \Setabs{x}{\cfind{e}(x) \fdefined}.
\]
Let $f$ be the function $f(x) = \tuple{e, x}$. Then for every natural
number $x$, $x \in B$ if and only if $f(x) \in K_0$. In other words, $f$
reduces $B$ to~$K_0$.

To see that $K_1$ is complete, note that in the proof of
\olref[k1]{prop:k1} we reduced $K_0$ to it. So, by
\olref[ppr]{prop:trans-red}, any computably enumerable set can be
reduced to~$K_1$ as well.

$K$ can be reduced to $K_0$ in much the same way.
\end{proof}

\begin{prob}
Give a reduction of $K$ to $K_0$.
\end{prob}

\begin{digress}
So, it turns out that all the examples of computably enumerable sets
that we have considered so far are either computable, or complete.
This should seem strange!{} Are there any examples of computably
enumerable sets that are neither computable nor complete? The answer
is yes, but it wasn't until the middle of the 1950s that this was
established by Friedberg and Muchnik, independently.
\end{digress}



% ==== OLP file: content/computability/computability-theory/computability-theory.tex ====
% Part: computability
% Chapter: computability-theory



\olchapter{cmp}{thy}{Computability Theory}

\begin{editorial}
  Material in this chapter should be reviewed and expanded. In
  particular, there are no exercises yet.
\end{editorial}

\olimport{introduction}

\olimport{coding-computations}

\olimport{normal-form}

\olimport{s-m-n}

\olimport{universal-part-function}

\olimport{no-universal-function}

\olimport{halting-problem}

\olimport{russells-paradox}

\olimport{computable-sets}

\olimport{ce-sets}

\olimport{equiv-ce-defs}

\olimport{non-comp-set}

\olimport{ce-closed-cup-cap}

\olimport{complement-ce}

\olimport{reducibility}

\olimport{prop-reduce}

\olimport{complete-ce-sets}

\olimport{k-1}

\olimport{total}

\olimport{rice-theorem}

\olimport{fixed-point-thm}

\olimport{application-fixed-point}

\olimport{def-functions-self-reference}

\OLEndChapterHook


% ==== OLP file: content/computability/computability-theory/computable-sets.tex ====
% Part: computability
% Chapter: computability-theory
% Section: computable-sets



\olfileid{cmp}{thy}{cps}
\olsection{Computable Sets}

We can extend the notion of computability from computable functions to
computable sets:

\begin{defn}
  Let $S$ be a set of natural numbers. Then $S$ is \emph{computable}
  iff its characteristic function~$\Char{S}$ is. In other words,
  $S$~is computable iff the function
\[
\Char{S}(x) =
\begin{cases}
1 & \text{if $x \in S$} \\
0 & \text{otherwise}
\end{cases}
\]
is computable. Similarly, a relation $R(x_0, \dots, x_{k-1})$ is
computable if and only if its characteristic function is.

Computable sets and relations are also called \emph{decidable}.
\end{defn}

\begin{explain}
Notice that we now have a number of notions of computability: for
partial functions, for functions, and for sets. Do not get them
confused!{} \iftag{TMs}{The Turing machine computing a partial function
  returns the output of the function, for input values at which the
  function is defined; the Turing machine computing a set returns
  either $1$ or~$0$, after deciding whether or not the input value is in
  the set or not.}{}
\end{explain}


% ==== OLP file: content/computability/computability-theory/def-functions-self-reference.tex ====
% Part: computability
% Chapter: computability-theory
% Section: def-functions-self-reference



\olfileid{cmp}{thy}{slf}
\olsection{Defining Functions using Self-Reference}

It is generally useful to be able to define functions in terms of
themselves. For example, given computable functions $k$, $l$, and~$m$,
the fixed-point lemma tells us that there is a partial computable
function~$f$ satisfying the following equation for every~$y$:
\[
f(y) \simeq
\begin{cases}
k(y) & \text{if $l(y) = 0$} \\
f(m(y)) & \text{otherwise.}
\end{cases}
\]
Again, more specifically, $f$~is obtained by letting
\[
g(x,y) \simeq
\begin{cases}
k(y) & \text{if $l(y) = 0$} \\
\cfind{x}(m(y)) & \text{otherwise}
\end{cases}
\]
and then using the fixed-point lemma to find an index~$e$ such that
$\cfind{e}(y) = g(e,y)$.

For a concrete example, the ``greatest common divisor'' function
$\fn{gcd}(u,v)$ can be defined by
\[
\fn{gcd}(u,v) \simeq
\begin{cases}
v & \text{if $u = 0$} \\
\fn{gcd}(\fn{mod}(v, u), u) & \text{otherwise}
\end{cases}
\]
where $\fn{mod}(v, u)$ denotes the remainder of dividing $v$
by~$u$. An appeal to the fixed-point lemma shows that $\fn{gcd}$ is
partial computable. (In fact, this can be put in the format above,
letting $y$ code the pair $\tuple{u, v}$.) A subsequent induction
on~$u$ then shows that, in fact, $\fn{gcd}$ is total.

Of course, one can cook up self-referential definitions that are much
fancier than the examples just discussed. Most programming languages
support definitions of functions in terms of themselves, one way or
another. Note that this is a little bit less dramatic than being able
to define a function in terms of an \emph{index} for an algorithm
computing the functions, which is what, in full generality, the
fixed-point theorem lets you do.


% ==== OLP file: content/computability/computability-theory/equiv-ce-defs.tex ====
% Part: computability
% Chapter: computability-theory
% Section: equiv-ce-defs



\olfileid{cmp}{thy}{eqc}

\olsection[Definitions of C. E. Sets]{Equivalent Defininitions of
  Computably Enumerable Sets}


The following gives a number of important equivalent statements of
what it means to be computably enumerable.

\begin{thm}
\ollabel{thm:ce-equiv}
Let $S$ be a set of natural numbers. Then the following are
equivalent:
\begin{enumerate}
\item\ollabel{case:ce} $S$ is computably enumerable.
\item\ollabel{case:ran-pc} $S$ is the range of a \emph{partial} computable function.
\item\ollabel{case:ran-prim} $S$ is empty or the range of a primitive recursive function.
\item\ollabel{case:ce-domain} $S$ is the \emph{domain} of a partial computable function.
\end{enumerate}
\end{thm}

\begin{explain}
The first three clauses say that we can equivalently take any non-empty
computably enumerable set to be enumerated by either a computable
function, a partial computable function, or a primitive recursive
function. The fourth clause tells us that if $S$ is computably
enumerable, then for some index~$e$,
\[
S = \Setabs{x}{\cfind{e}(x) \fdefined}.
\]
In other words, $S$ is the set of inputs on for which the computation
of $\cfind{e}$ halts. For that reason, computably enumerable sets are
sometimes called \emph{semi-decidable}: if a number is in the set, you
eventually get a ``yes,'' but if it isn't, you never get a ``no''!{}
\end{explain}

\begin{proof}
Since every primitive recursive function is computable and every
computable function is partial computable, \olref{case:ran-prim}
implies \olref{case:ce} and \olref{case:ce}
implies~\olref{case:ran-pc}. (Note that if $S$ is empty, $S$~is the
range of the partial computable function that is nowhere defined.) If
we show that \olref{case:ran-pc} implies \olref{case:ran-prim}, we
will have shown the first three clauses equivalent.

So, suppose $S$ is the range of the partial computable function
$\cfind{e}$. If $S$ is empty, we are done. Otherwise, let $a$ be any
element of~$S$. By Kleene's normal form theorem, we can write
\[
\cfind{e}(x) = U(\umin{s}{T(e, x, s)}).
\]
In particular, $\cfind{e}(x) \fdefined$ and $= y$ if and only if there
is an $s$ such that $T(e, x, s)$ and $U(s) = y$. Define $f(z)$ by
\[
f(z) = \begin{cases}
  U((z)_1) & \text{if $T(e, (z)_0, (z)_1)$} \\
  a        & \text{otherwise.}
\end{cases}
\]
Then $f$ is primitive recursive, because $T$ and $U$
are. \iftag{TMs}{Expressed in terms of Turing machines, if $z$ codes a
  pair $\tuple{(z)_0, (z)_1}$ such that $(z)_1$ is a halting
  computation of machine~$M_e$ on input $(z)_0$, then $f$ returns the
  output of the computation; otherwise, it returns~$a$.}

We need to show that $S$ is the range of~$f$, i.e., for any natural
number~$y$, $y \in S$ if and only if it is in the range of~$f$. In the
forwards direction, suppose $y \in S$. Then $y$ is in the range of
$\cfind{e}$, so for some $x$ and~$s$, $T(e,x,s)$ holds and $U(s) = y$.
But then $y = f(\tuple{x,s})$. Conversely, suppose $y$ is in the range
of~$f$. Then either $y = a$, or for some~$z$, $T(e,(z)_0,(z)_1)$ and
$U((z)_1) = y$. Since, in the latter case, $\cfind{e}(x) \fdefined =
y$, either way, $y$ is in~$S$.

(The notation $\cfind{e}(x) \fdefined = y$ means ``$\cfind{e}(x)$ is
defined and equal to $y$.'' We could just as well use $\cfind{e}(x) =
y$, but the extra arrow is sometimes helpful in reminding us that we
are dealing with a partial function.)

To finish up the proof of \olref{thm:ce-equiv}, it suffices to show
that \olref{case:ce} and~\olref{case:ce-domain} are equivalent. First,
let us show that \olref{case:ce} implies~\olref{case:ce-domain}.
Suppose $S$ is the range of a computable function~$f$, i.e.,
\[
S = \Setabs{y}{\text{for some $x$, } f(x) = y}.
\]
Let
\[
g(y) = \umin{x}{(f(x) = y)}.
\]
Then $g$ is a partial computable function, and $g(y)$ is defined if
and only if for some~$x$, $f(x) = y$. In other words, the domain of
$g$~is the range of~$f$. \iftag{TMs}{Expressed in terms of Turing
machines: given a Turing machine~$F$ that enumerates the elements
of~$S$, let $G$ be the Turing machine that semi-decides $S$ by
searching through the outputs of~$F$ to see if a given element is in
the set, halts if it is and keeps searching forever if it isn't.}{}

Finally, to show \olref{case:ce-domain} implies~\olref{case:ce},
suppose that $S$~is the domain of the partial computable
function~$\cfind{e}$, i.e.,
\[
S = \Setabs{x}{\cfind{e}(x) \fdefined}.
\]
If $S$ is empty, we are done; otherwise, let $a$ be any element
of~$S$. Define $f$ by
\[
f(z) = \begin{cases}
(z)_0 & \text{if $T(e,(z)_0,(z)_1)$} \\
a & \text{otherwise.}
\end{cases}
\]
Then, as above, a number $x$ is in the range of~$f$ if and only if
$\cfind{e}(x) \fdefined$, i.e., if and only if $x \in S$. \iftag{TMs}{Expressed
in terms of Turing machines: given a machine $M_e$ that semi-decides
$S$, enumerate the elements of $S$ by running through all possible
Turing machine computations, and returning the inputs that correspond
to halting computations.}{}
\end{proof}

Clause~\olref{case:ce-domain} of \olref{thm:ce-equiv} provides us with
a convenient way of enumerating the computably enumerable sets: for
each~$e$, let $W_e$ denote the domain of $\cfind{e}$, i.e.,
\[
W_e = \Setabs{x}{\cfind{e}(x) \fdefined}.
\] 
Then if $A$ is any computably enumerable set, $A = W_e$, for some~$e$.

The following provides yet another characterization of the computably
enumerable sets.

\begin{thm}
\ollabel{thm:exists-char}
A set $S$ is computably enumerable if and only if there is a
computable relation $R(x,y)$ such that
\[
S = \Setabs{ x }{ \lexists[y][R(x,y)] }.
\]
\end{thm}

\begin{proof}
In the forward direction, suppose $S$ is computably
enumerable. Then for some $e$, $S = W_e$. For this value of~$e$
we can write $S$ as
\[
S = \Setabs{ x }{ \lexists[y][T(e, x, y)] }.
\]
In the reverse direction, suppose $S = \Setabs{ x }{
  \lexists[y][R(x, y)] }$. Define $f$~by
\[
f(x) \simeq \umin{y}{R(x, y)}.
\]
Then $f$ is partial computable, and $S$ is the domain of~$f$.
\end{proof}



% ==== OLP file: content/computability/computability-theory/fixed-point-thm.tex ====
% Part: computability
% Chapter: computability-theory
% Section: fixed-point-thm



\olfileid{cmp}{thy}{fix}
\olsection{The Fixed-Point Theorem}

Let's consider the halting problem again. As temporary
notation, let us write $\gn{\cfind{x}(y)}$ for $\tuple{x, y}$; think of
this as representing a ``name'' for the value $\cfind{x}(y)$. With this
notation, we can reword one of our proofs that the halting problem is
undecidable.

Question: is there a computable function $h$, with the
following property? For every $x$ and $y$,
\[
h(\gn{\cfind{x}(y)}) =
\begin{cases}
1 & \text{if $\cfind{x}(y) \fdefined$} \\
0 & \text{otherwise.}
\end{cases}
\]

Answer: No; otherwise, the partial function
\[
g(x) \simeq
\begin{cases}
0 & \text{if $h(\gn{\cfind{x}(x)}) = 0$} \\
\text{undefined} & \text{otherwise}
\end{cases}
\]
would be computable, and so have some index~$e$. But then we have
\[
\cfind{e}(e) \simeq
\begin{cases}
0 & \text{if $h(\gn{\cfind{e}(e)}) = 0$} \\
\text{undefined} & \text{otherwise,}
\end{cases}
\]
in which case $\cfind{e}(e)$ is defined if and only if it isn't, a
contradiction.

Now, take a look at the equation with $\cfind{e}$. There is an instance of
self-reference there, in a sense: we have arranged for the value of
$\cfind{e}(e)$ to depend on $\gn{\cfind{e}(e)}$, in a certain way. The
fixed-point theorem says that we {\em can} do this, in general---not
just for the sake of proving contradictions.

\olref{lem:fixed-equiv} gives two equivalent ways of stating the
fixed-point theorem. Logically speaking, the fact that the statements
are equivalent follows from the fact that they are both true; but what
we really mean is that each one follows straightforwardly from the
other, so that they can be taken as alternative statements of the same
theorem.

\begin{lem}
\ollabel{lem:fixed-equiv}
The following statements are equivalent:
\begin{enumerate}
\item For every partial computable function $g(x,y)$, there is an
  index~$e$ such that for every~$y$,
\[
\cfind{e}(y) \simeq g(e,y).
\]
\item For every computable function~$f(x)$, there is an index~$e$ such
  that for every~$y$,
\[
\cfind{e}(y) \simeq \cfind{f(e)}(y).
\]
\end{enumerate}
\end{lem}

\begin{proof}
$(1) \Rightarrow (2)$: Given $f$, define $g$ by $g(x,y) \simeq
\fn{Un}(f(x),y)$. Use (1) to get an index~$e$ such that for every~$y$,
\begin{align*}
\cfind{e}(y) & = \fn{Un}(f(e),y) \\
& = \cfind{f(e)}(y).
\end{align*}

$(2) \Rightarrow (1)$: Given $g$, use the $s$-$m$-$n$ theorem to get $f$ such
that for every $x$ and~$y$, $\cfind{f(x)}(y) \simeq g(x,y)$. Use (2) to
get an index~$e$ such that
\begin{align*}
\cfind{e}(y) & = \cfind{f(e)}(y) \\
& = g(e,y).
\end{align*}
This concludes the proof.
\end{proof}

\begin{explain}
Before showing that statement (1) is true (and hence (2) as well),
consider how bizarre it is. Think of $e$ as being a computer program;
statement (1) says that given any partial computable $g(x,y)$, you can
find a computer program $e$ that computes $g_e(y) \simeq g(e,y)$. In
other words, you can find a computer program that computes a function
that references the program itself.
\end{explain}

\begin{thm}
The two statements in \olref{lem:fixed-equiv} are
true. Specifically, for every partial computable function $g(x,y)$,
there is an index~$e$ such that for every~$y$,
\[
\cfind{e}(y) \simeq g(e,y).
\]
\end{thm}

\begin{proof}
The ingredients are already implicit in the discussion of the halting
problem above. Let $\fn{diag}(x)$ be a computable function which for each
$x$ returns an index for the function $f_x(y) \simeq \cfind{x}(x,y)$,
i.e.
\[
\cfind{\fn{diag}(x)}(y) \simeq \cfind{x}(x,y).
\]
Think of $\fn{diag}$ as a function that transforms a program for a 2-ary
function into a program for a 1-ary function, obtained by fixing the
original program as its first argument. The function $\fn{diag}$ can be
defined formally as follows: first define $s$ by
\[
s(x,y) \simeq \fn{Un}^2(x,x,y),
\]
where $\fn{Un}^2$ is a 3-ary function that is universal for partial computable
2-ary functions. Then, by the $s$-$m$-$n$ theorem, we can find a primitive
recursive function $\fn{diag}$ satisfying
\[
\cfind{\fn{diag}(x)}(y) \simeq s(x,y).
\]

Now, define the function $l$ by
\[
l(x,y) \simeq g(\fn{diag}(x),y).
\]
and let $\gn{l}$ be an index for $l$. Finally, let $e = \fn{diag}(\gn{l})$.
Then for every $y$, we have
\begin{align*}
\cfind{e}(y) & \simeq \cfind{\fn{diag}(\gn{l})}(y) \\
& \simeq \cfind{\gn{l}}(\gn{l}, y) \\
& \simeq l(\gn{l}, y) \\
& \simeq g(\fn{diag}(\gn{l}),y) \\
& \simeq g(e, y),
\end{align*}
as required.
\end{proof}

\begin{explain}
What's going on? Suppose you are given the task of writing a computer
program that prints itself out. Suppose further, however, that you are
working with a programming language with a rich and bizarre library of
string functions. In particular, suppose your programming language has
a function $\fn{diag}$ which works as follows: given an input
string~$s$, $\fn{diag}$ locates each instance of the symbol `x'
occurring in~$s$, and replaces it by a quoted version of the original
string. For example, given the string
\begin{quote}
\begin{verbatim}
hello x world
\end{verbatim}
\end{quote}
as input, the function returns
\begin{quote}
\begin{verbatim}
hello 'hello x world' world
\end{verbatim}
\end{quote}
as output. In that case, it is easy to write the desired program; you
can check that
\begin{quote}
\begin{verbatim}
print(diag('print(diag(x))'))
\end{verbatim}
\end{quote}
does the trick. For more common programming languages like C++ and
Java, the same idea (with a more involved implementation) still works.

We are only a couple of steps away from the proof of the fixed-point
theorem. Suppose a variant of the print function $\fn{print}(x,y)$
accepts a string $x$ and another numeric argument $y$, and prints the
string $x$ repeatedly, $y$ times. Then the ``program''
\begin{quote}
\begin{verbatim}
getinput(y);
print(diag('getinput(y); print(diag(x), y)'), y)
\end{verbatim}
\end{quote}
prints itself out $y$ times, on input $y$. Replacing the
$\fn{getinput}$---$\fn{print}$---$\fn{diag}$ skeleton by an
arbitrary function $g(x,y)$ yields
\begin{quote}
\begin{verbatim}
g(diag('g(diag(x), y)'), y)
\end{verbatim}
\end{quote}
which is a program that, on input $y$, runs $g$ on the program itself
and $y$. Thinking of ``quoting'' with ``using an index for,'' we have
the proof above.

For now, it is o.k.\ if you want to think of the proof as formal
trickery, or black magic. But you should be able to reconstruct the
details of the argument given above. When we prove the incompleteness
theorems (and the related ``fixed-point theorem'') we will discuss
other ways of understanding why it works.
\end{explain}

\begin{tagblock}{lambda}
\begin{digress}
The same idea can be used to get a ``fixed point'' combinator. Suppose
you have a lambda term $g$, and you want another term $k$ with the
property that $k$ is $\beta$-equivalent to $gk$. Define terms
\[
\fn{diag}(x) = xx
\]
and
\[
l(x) = g(\fn{diag}(x))
\]
using our notational conventions; in other words, $l$ is the term
$\lambd[x][g(xx)]$. Let $k$ be the term $ll$. Then we have
\begin{align*}
k & = (\lambd[x][g(xx)])(\lambd[x][g(xx)]) \\
& \red  g((\lambd[x][g(xx)])(\lambd[x][g(xx)])) \\
& = gk.
\end{align*}
If one takes
\[
Y = \lambd[g][((\lambd[x][g(xx)])(\lambd[x][g(xx)]))]
\]
then $Yg$ and $g(Yg)$ reduce to a common term; so $Yg \equiv_\beta
g(Yg)$. This is known as ``Curry's combinator.'' If instead one takes
\[
Y = (\lambd[xg][g(xxg)])(\lambd[xg][g(xxg)])
\]
then in fact $Yg$ reduces to $g(Yg)$, which is a stronger statement.
This latter version of $Y$ is known as ``Turing's combinator.''
\end{digress}
\end{tagblock}


% ==== OLP file: content/computability/computability-theory/halting-problem.tex ====
% Part: computability
% Chapter: computability-theory
% Section: halting-problem



\olfileid{cmp}{thy}{hlt}
\olsection{The Halting Problem}

By construction, the universal partial computable
function~$\fn{Un}(e,x)$ is defined if and only if the computation of
the function coded by~$e$ produces a value for input~$x$. It is
natural to ask if we can decide whether this is the case. In fact, it
is not. For the Turing machine model of computation, this means that
whether a given Turing machine halts on a given input is
computationally undecidable. The following theorem is therefore known
as the ``undecidability of the halting problem.'' We will provide two
proofs below. The first continues the thread of our previous
discussion, while the second is more direct.

\begin{thm}
\ollabel{thm:halting-problem}
Let
\[
h(e, x)  =
\begin{cases}
1 & \text{if\/ $\fn{Un}(e, x)$ is defined} \\
0 & \text{otherwise.}
\end{cases}
\]
Then $h$ is not computable.
\end{thm}

\begin{proof}
Suppose $h$ is computable. We show that this would let us define a
universal computable function. Define
\[
\fn{Un'}(e,x) =
\begin{cases}
\fn{Un}(e,x) & \text{if $h(e,x) = 1$} \\
0 & \text{otherwise.}
\end{cases}
\]
But now $\fn{Un'}(e, x)$ is a total function, and is computable if~$h$
is. For instance, we could define $g$ using primitive recursion, by
\begin{align*}
g(0, e, x) & \simeq 0 \\
g(y+1, e, x) & \simeq \fn{Un}(e,x);
\end{align*}
then
\[
\fn{Un'}(e,x) \simeq g(h(e,x),e,x).
\]
Since $\fn{Un'}(e,x)$ agrees with $\fn{Un}(e,x)$ wherever the latter
is defined, $\fn{Un'}$ is universal for those partial computable
functions that happen to be total. But this contradicts
\olref[nou]{thm:no-univ}.
\end{proof}

\begin{proof}
Suppose $h(e,x)$ were computable. Define the function $g$ by
\[
g(x) =
\begin{cases}
  0                & \text{if $h(x,x) = 0$} \\
  \fundefined & \text{otherwise.}
\end{cases}
\]
The function $g$ is partial computable. For example, one can define it
as $\umin{y}{h(x,x) = 0}$. So, for some~$e$, $g(x) \simeq \fn{Un}(e,
x)$ for every~$x$. Is $g$ defined at $e$?  If it is, then, by the
definition of~$g$, $h(e,e) = 0$ ($h$~can only take the value~$0$ if it
is defined). By the definition of~$h$, this means that $\fn{Un}(e, e)$
is undefined. By our assumption that $g(x) \simeq \fn{Un}(e, x)$
for every~$x$, we have that $g(e)$ is undefined, a contradiction.
On the other hand, if $g(e)$ is undefined, then $h(e,e) \neq 0$, and
so $h(e,e) = 1$. It follows that $\fn{Un}(e, e)$ is defined. But since
$g(x) \simeq \fn{Un}(e, x)$, then $g(e)$ would also be defined. Again,
a contradiction.
\end{proof}

\begin{tagblock}{TMs}
\begin{explain}
We can describe this argument in terms of Turing machines.  Suppose
there were a Turing machine~$H$ that takes as input a description of a
Turing machine~$E$ and an input~$x$, and decides whether or not $E$
halts on input~$x$. Then we could build another Turing machine~$G$
which takes a single input~$x$, runs $H$ to decide if the machine
$M_x$ with index~$x$ halts on input~$x$, and does the opposite. In
other words, if $H$ reports that $M_x$ halts on input~$x$, $G$ goes
into an infinite loop, and if $H$ reports that $M_x$ doesn't halt on
input~$x$, then $G$ just halts. Does $G$ halt on its own index as
input? The argument above shows that it does if and only if it
doesn't---a contradiction. So our supposition that there is a such
Turing machine~$H$ must be false.
\end{explain}
\end{tagblock}


% ==== OLP file: content/computability/computability-theory/introduction.tex ====
% Part: computability
% Chapter: computability-theory
% Section: introduction



\olfileid{cmp}{thy}{int}
\olsection{Introduction}

The branch of logic known as \emph{computability theory} deals with
issues having to do with the computability, or relative computability,
of functions and sets. It is evidence of Kleene's influence
that the subject used to be known as \emph{recursion theory}, and
today, both names are commonly used.

Let us call a function~$f\colon \Nat \pto \Nat$ \emph{partial
  computable} if it can be computed in some model of computation. If
$f$ is total we will simply say that $f$~is \emph{computable}. A
relation~$R$ with computable characteristic function~$\Char{R}$ is
also called computable. If $f$ and~$g$ are partial functions, we will
write $f(x) \fdefined$ to mean that $f$ is defined at~$x$, i.e., $x$~is
in the domain of~$f$; and $f(x) \fundefined$ to mean the opposite,
i.e., that $f$ is not defined at~$x$. We will use $f(x) \simeq g(x)$
to mean that either $f(x)$ and $g(x)$ are both undefined, or they are
both defined and equal.

One can explore the theory of computability without having to refer to
a specific model of computation. To do this, one shows that there is a
universal partial computable function~$\fn{Un}(k, x)$. This allows us
to enumerate the partial computable functions. We will adopt the
notation~$\cfind{k}$ to denote the $k$-th unary partial computable
function, defined by $\cfind{k}(x) \simeq \fn{Un}(k, x)$. (Kleene used
$\{ k \}$ for this purpose, but this notation has not been used as
much recently.)  Slightly more generally, we can uniformly enumerate
the partial computable functions of arbitrary arities, and we will use
$\cfind{k}[n]$ to denote the $k$-th $n$-ary partial recursive
function.

If $f(\vec x, y)$ is a total or partial function, then $\umin{y}{f
(\vec x, y)}$ is the function of~$\vec x$ that returns the least~$y$
such that $f(\vec x, y) = 0$, assuming that all of $f(\vec x, 0)$,
\dots, $f(\vec x, y-1)$ are defined; if there is no such $y$,
$\umin{y}{f (\vec x, y)}$ is undefined. If $R(\vec x, y)$ is a
relation, $\umin{y}{R(\vec x, y)}$ is defined to be the least~$y$ such
that $R(\vec x, y)$ is true; in other words, the least~$y$ such that
$1 \tsub \Char{R}(\vec x, y) = 0$.

To show that a function is computable, there are
two ways one can proceed:
\begin{enumerate}
\item Rigorously: describe a Turing machine or partial recursive
  function explicitly, and show that it computes the function you have
  in mind;
\item Informally: describe an algorithm that computes it, and appeal to
  Church's thesis.
\end{enumerate}
There is no fine line between the two; a detailed description of
an algorithm should provide enough information so that it is
relatively clear how one could, in principle, design the right Turing
machine or sequence of partial recursive definitions. Fully rigorous
definitions are unlikely to be informative, and we will try to find a
happy medium between these two approaches; in short, we will try to
find intuitive yet rigorous proofs that the precise definitions could
be obtained.


% ==== OLP file: content/computability/computability-theory/k-1.tex ====
% Part: computability
% Chapter: computability-theory
% Section: k-1



\olfileid{cmp}{thy}{k1}
\olsection{An Example of Reducibility}

Let us consider an application of \olref[ppr]{prop:reduce}.

\begin{prop}
\ollabel{prop:k1}
Let
\[
K_1 = \Setabs{e}{\cfind{e}(0) \fdefined}.
\]
Then $K_1$ is computably enumerable but not computable.
\end{prop}

\begin{proof}
Since $K_1 = \Setabs{e}{\lexists[s][T(e,0,s)]}$, $K_1$ is computably
enumerable by \olref[eqc]{thm:exists-char}.

To show that $K_1$ is not computable, let us show that $K_0$ is
reducible to it.

\begin{explain}
This is a little bit tricky, since using $K_1$ we can
only ask questions about computations that start with a particular
input, $0$. Suppose you have a smart friend who can answer questions
of this type (friends like this are known as ``oracles''). Then
suppose someone comes up to you and asks you whether or not $\tuple{e,
  x}$ is in $K_0$, that is, whether or not machine $e$ halts on input
$x$. One thing you can do is build another machine, $e_x$, that, for
\emph{any} input, ignores that input and instead runs~$e$ on input
$x$. Then clearly the question as to whether machine $e$ halts on
input $x$ is equivalent to the question as to whether machine $e_x$
halts on input $0$ (or any other input). So, then you ask your friend
whether this new machine, $e_x$, halts on input $0$; your friend's
answer to the modified question provides the answer to the original
one. This provides the desired reduction of $K_0$ to~$K_1$.
\end{explain}

Using the universal partial computable function, let $f$
be the 3-ary function defined by
\[
f(x,y,z) \simeq \cfind{x}(y).
\]
Note that $f$ ignores its third input entirely. Pick an index $e$ such
that $f = \cfind{e}[3]$; so we have
\[
\cfind{e}[3](x,y,z) \simeq \cfind{x}(y).
\]
By the $s$-$m$-$n$ theorem, there is a function $s(e,x,y)$ such that, for
every $z$,
\begin{align*}
\cfind{s(e,x,y)}(z) & \simeq \cfind{e}[3](x,y,z) \\
& \simeq \cfind{x}(y).
\end{align*}

\begin{explain}
In terms of the informal argument above, $s(e,x,y)$ is an index for
the machine that, for any input $z$, ignores that input and computes
$\cfind{x}(y)$.
\end{explain}

In particular, we have
\[
\cfind{s(e,x,y)}(0) \fdefined \quad \text{if and only if} \quad
\cfind{x}(y) \fdefined.
\]
In other words, $\tuple{x, y} \in K_0$ if and only if $s(e,x,y) \in
K_1$. So the function $g$ defined by
\[
g(w) = s(e,(w)_0,(w)_1)
\]
is a reduction of $K_0$ to~$K_1$.
\end{proof}


% ==== OLP file: content/computability/computability-theory/no-universal-function.tex ====
% Part: computability
% Chapter: computability-theory
% Section: no-universal-function



\olfileid{cmp}{thy}{nou}
\olsection{No Universal Computable Function}

Although there is a partial computable function that is total for the
partial computable functions, there is no total computable
function that is universal for the total computable functions. 

\begin{thm}
\ollabel{thm:no-univ}
There is no universal computable function. In other words, any
function $\fn{Un}'(k, x)$ which is such that if $f(x)$ is a total
computable function, then there is a natural number~$k$ such that
$f(x) = \fn{Un}'(k,x)$ for every~$x$, is not computable.
\end{thm}

\begin{proof}
The proof is a simple diagonalization: if $\fn{Un}'(k,x)$ were total
and computable, then
\[
d(x) = \fn{Un}'(x, x) + 1
\]
would also be total and computable. However, by definition, $d(k)$ is
not equal to $\fn{Un}'(k,k)$. Hence, for every $k$, the values of
$d(x)$ and~$\fn{Un}'(k, x)$ differ for at least one~$x$, namely $x = k$.
\end{proof}

\begin{explain}
\olref[uni]{thm:univ-comp} above shows that we can get around this
diagonalization argument, but only at the expense of allowing the
universal function to be partial. That is, $\fn{Un}$ is universal for
the total computable functions, it just isn't total. The
diagonalization argument doesn't work in the partial case. 
\end{explain}

\begin{prob}
  To understand why the diagonalization argument in the proof of
  \olref{thm:no-univ} does not work in the partial
  case, consider the function $f(x) \simeq \fn{Un}(x,x)+1$. Is it
  partial computable? If so, it has an index~$e$, i.e., $f(x) \simeq
  \fn{Un}(e,x)$. What can you say about~$f(e)$?
\end{prob}



% ==== OLP file: content/computability/computability-theory/non-comp-set.tex ====
% Part: computability
% Chapter: computability-theory
% Section: non-comp-set



\olfileid{cmp}{thy}{ncp}
\olsection{There Are Non-Computable Sets}


We saw above that every computable set is computably enumerable. Is
the converse true? The following shows that, in general, it is not.

\begin{thm}\ollabel{thm:K-0}
Let $K_0$ be the set $\Setabs{\tuple{e, x}}{\cfind{e}(x) \fdefined}$.
Then $K_0$ is computably enumerable but not computable.
\end{thm}

\begin{proof}
To see that $K_0$ is computably enumerable, note that it is the
domain of the function~$f$ defined by
\[
f(z) = \umin{y}{(\len{z} = 2 \land T((z)_0, (z)_1, y))}.
\]
For, if $\cfind{e}(x)$ is defined, $f(\tuple{e, x})$ finds a halting
computation sequence; if $\cfind{e}(x)$ is undefined, so is
$f(\tuple{e, x})$; and if $z$ doesn't even code a pair, then $f(z)$ is
also undefined.

The fact that $K_0$ is not computable is just the undecidability of
the halting problem, \olref[hlt]{thm:halting-problem}.
\end{proof}

The set $K_0$ is the set of pairs $\tuple{e,x}$ such that
$\cfind{e}(x) \fdefined$, i.e., $\tuple{e,x} \in K_0$ iff $\cfind{e}$
is defined (halts) on input~$x$, so it is also called the ``halting
set.'' The set $K = \Setabs{e}{\cfind{e}(e) \fdefined}$ is the
``self-halting set.'' It is often used as a canonical undecidable set.

\begin{thm}\ollabel{thm:K}
The self-halting set $K = \Setabs{e}{\cfind{e}(e) \fdefined}$ is
!!{c.e.} but not decidable.
\end{thm}

\begin{proof}
  Suppose $K$ is decidable, i.e., its characteristic function
  $\Char{K}$ is computable. Let 
  \[d(e) = \begin{cases}
  1 & \text{if\/ $\Char{K}(e) = 0$}\\
  \fundefined & \text{otherwise.}
  \end{cases}
  \] 
  Let $k$ be the index of~$d$, i.e., $d \simeq \cfind{k}$. Then $d(k)
  \simeq \cfind{k}(k)$. This contradicts the fact that $d(k)
  \fdefined$ iff $\cfind{k}(k) \fundefined$, which follows from the
  definition of~$d$.

  $K$ is the domain of $f(x) = \umin{y}{T(x,x,y)}$ and so is !!{c.e.}
\end{proof}


% ==== OLP file: content/computability/computability-theory/normal-form.tex ====
% Part: computability
% Chapter: computability-theory
% Section: normal-form



\olfileid{cmp}{thy}{nfm}
\olsection{The Normal Form Theorem}

Suppose we can describe definitions of computable functions, and test if
some putative description of the complete record of the computation of
that function on some input is correct. Then it stands to reason that
independently of the model of computation, we can determine the value
of any computable function~$f$ on any input~$x$ as follows:
\begin{enumerate}
  \item Search through all possible descriptions of records of
  computation.
  \item Test if a given record is the record of a computation of~$f(x)$.
  \item Extract the value of~$f(x)$ from the correct record if we have
  found it.
\end{enumerate}
That this is in fact true is the content of Kleene's normal form
theorem.

\begin{thm}[Kleene's Normal Form Theorem]
\ollabel{thm:normal-form}
There is a primitive recursive relation~$T(e, x, s)$ and a primitive
recursive function~$U(s)$, with the following property: if $f$ is any
partial computable function, then for some~$e$,
\[
f(x) \simeq U(\umin{s}{T(e, x, s)})
\]
for every~$x$.
\end{thm}

\begin{proof}[Proof Sketch]
For any model of computation one can rigorously define a description
of the computable function~$f$ and code such description using a
natural number~$e$.  One can also rigorously define a notion of
``computation sequence'' which records the process of computing the
function with index~$e$ for input~$x$.  Such a computation sequence can
likewise be coded as a number~$s$.  This can be done in such a way
that
\begin{enumerate}
  \item the relation $T(e, x, s)$, which holds iff a number~$s$ codes the computation
  sequence of the function with index~$e$ on input~$x$, and
  \item the function $U(s)$ which maps a computation sequence coded
  by~$s$ to the end result of that computation
\end{enumerate}
are both computable. In fact, the relation~$T$ and the function~$U$
are primitive recursive.
\end{proof}

\begin{explain}
In order to give a rigorous proof of the Normal Form Theorem, we would
have to fix a model of computation and carry out the coding of
descriptions of computable functions and of computation sequences in
detail, and verify that the relation~$T$ and function~$U$ are
primitive recursive.  For most applications, it suffices that $T$
and~$U$ are computable and that $U$~is total.

It is probably best to remember the proof of the normal form theorem
in slogan form: $\umin{s}{T(e, x, s)}$ searches for a computation
sequence of the function with index~$e$ on input~$x$, and $U$ returns
the output of the computation sequence if one can be found.
\end{explain}

If the model of computation is the partial recursive functions (which
is what Kleene originally used), it shows that only a single use of
unbounded search, i.e., a single $\umin{y}{f(\vec x, y)}$ operator is
necessary for the definition of any function. In this sense it shows
that any partial recursive function has a normal form.

$T$ and $U$ can be used to define the enumeration $\cfind{0}$,
$\cfind{1}$, $\cfind{2}$, \dots.  From now on, we will assume that we
have fixed a suitable choice of $T$ and~$U$, and take the equation
\[
\cfind{e}(x) \simeq U(\umin{s}{T(e,x,s)})
\]
to be the \emph{definition} of $\cfind{e}$.

Here is another useful fact:

\begin{thm}
Every partial computable function has infinitely many indices.
\end{thm}

Again, this is intuitively clear.  Given any (description of) a
computable function, one can come up with a different description
which computes the same function (input-output pair) but does so,
e.g., by first doing something that has no effect on the computation
(say, test if $0 = 0$, or count to $5$, etc.).  The index of the
altered description will always be different from the original
index.  Both are indices of the same function, just computed slightly
differently.


% ==== OLP file: content/computability/computability-theory/prop-reduce.tex ====
% Part: computability
% Chapter: computability-theory
% Section: prop-reduce



\olfileid{cmp}{thy}{ppr}
\olsection{Properties of Reducibility}

We write $A \leq_m B$ if $A$ reduces to~$B$, and this notation
suggests that if $A \leq_m B$, then $A$ is ``no harder than''~$B$ and
that $B$ is ``as hard or harder than''~$A$, and that $\le_m$, like the
usual $\le$ on numbers, orders sets (or decision problems) by their
complexity. The following two propositions support this intuition. The
first one say taht $\le_m$ is transitive.

\begin{prop}
\ollabel{prop:trans-red}
If $A \leq_m B$ and $B \leq_m C$, then $A \leq_m C$.
\end{prop}

\begin{proof}
Composing a reduction of $A$ to~$B$ with a reduction of $B$ to
$C$ yields a reduction of $A$ to~$C$.
\end{proof}

\begin{prob}
Prove \olref{prop:trans-red} by showing that if $f$ and~$g$ are
many-one reductions of $A$ to~$B$ and $B$ to~$C$, respectively, then
$\comp{f}{g}$ is a many-one reduction of $A$ to~$C$.
\end{prob}

\begin{prop}
\ollabel{prop:reduce}
Let $A$ and $B$ be any sets, and suppose $A \leq_m B$.
\begin{enumerate}
\item If $B$ is computably enumerable, so is~$A$.
\item If $B$ is computable, so is~$A$.
\end{enumerate}
\end{prop}

\begin{proof}
Let $f$ be a many-one reduction from $A$ to~$B$. For the first
claim, just check that if $B$ is the domain of a partial function~$g$,
then $A$ is the domain of~$\comp{f}{g}$:
\begin{align*}
x \in A & \text{ iff } f(x) \in B \\
& \text{ iff }  g(f(x)) \fdefined.
\end{align*}

For the second claim, remember that if $B$~is computable then $B$
and~$\Complement{B}$ are computably enumerable
(\olref[cmp]{thm:ce-comp}). It is not hard to check that $f$~is also a
many-one reduction of $\Complement{A}$ to $\Complement{B}$, so, by the
first part of this proof, $A$ and~$\Complement{A}$ are !!{computably
enumerable}. So $A$ is computable as well by \olref[cmp]{thm:ce-comp}.
(Alternatively, you can check that $\Char{A} = \comp{f}{\Char{B}}$; so
if $\Char{B}$ is computable, then so is~$\Char{A}$.)
\end{proof}

\begin{prob}
Suppose $f$ is a many-one reduction of $A$ to~$B$. Show that $f$ is
also a many-one reduction of $\Complement{A}$ to $\Complement{B}$.
\end{prob}

\begin{prob}
Show that if $f\colon A \to B$ is a many-one reduction, then $\Char{A}
= \comp{f}{\Char{B}}$.
\end{prob}

\begin{digress}
A more general notion of reducibility called \emph{Turing
reducibility} is useful in other contexts, especially for proving
undecidability results. Note that by \olref[cmp]{cor:comp-k}, the
complement of~$K_0$ is not reducible to~$K_0$, since it is not
computably enumerable. But, intuitively, if you knew the answers to
questions about $K_0$, you would know the answer to questions about
its complement as well. A set $A$ is said to be Turing reducible
to~$B$ if one can determine answers to questions in~$A$ using a
computable procedure that can ask questions about~$B$. This is more
liberal than many-one reducibility, in which (1)~you are only allowed
to ask one question about $B$, and (2)~a ``yes'' answer has to
translate to a ``yes'' answer to the question about $A$, and similarly
for ``no.'' It is still the case that if $A$~is Turing reducible
to~$B$ and $B$~is computable then $A$~is computable as well (though,
as we have seen, the analogous statement does not hold for computable
enumerability).

You should think about the various notions of reducibility we have
discussed, and understand the distinctions between them. We will,
however, only deal with many-one reducibility in this chapter.
Incidentally, both types of reducibility discussed in the last
paragraph have analogues in computational complexity, with the added
requirement that the Turing machines run in polynomial time: the
complexity version of many-one reducibility is known as \emph{Karp
reducibility}, while the complexity version of Turing reducibility is
known as \emph{Cook reducibility}.
\end{digress}


% ==== OLP file: content/computability/computability-theory/reducibility.tex ====
% Part: computability
% Chapter: computability-theory
% Section: reducibility



\olfileid{cmp}{thy}{red}
\olsection{Reducibility}

\begin{explain}
We now know that there is at least one set, $K_0$, that is computably
enumerable but not computable. It should be clear that there are
others. The method of reducibility provides a powerful method of
showing that other sets have these properties, without constantly
having to return to first principles.

Generally speaking, a ``reduction'' of a set $A$ to a set~$B$ is a
method of transforming answers to whether or not !!{element}s are in~$B$
into answers as to whether or not !!{element}s are in~$A$. We will focus
on a notion called ``many-one reducibility,'' but there are many other
notions of reducibility available, with varying properties. Notions of
reducibility are also central to the study of computational
complexity, where efficiency issues have to be considered as well. For
example, a set is said to be ``NP-complete'' if it is in NP and every
NP problem can be reduced to it, using a notion of reduction that is
similar to the one described below, only with the added requirement
that the reduction can be computed in polynomial time.

We have already used the notion of reduction notion implicitly. Define
the set~$K$ by
\[
K = \Setabs{x}{\cfind{x}(x) \fdefined},
\]
i.e., $K = \Setabs{x}{x \in W_x}$. Our proof that the halting problem
in unsolvable (\olref[hlt]{thm:halting-problem}) shows most directly
that $K$ is not computable. Recall that $K_0$ is the set
\[
K_0 = \Setabs{\tuple{e, x}}{\cfind{e}(x) \fdefined },
\]
i.e., $K_0 = \Setabs{\tuple{x,e}}{x \in W_e}$. It is easy to extend
any proof of the uncomputability of~$K$ to the uncomputability
of~$K_0$: if $K_0$ were computable, we could decide whether or not
!!a{element}~$x$ is in $K$ simply by asking whether or not the pair
$\tuple{x, x}$ is in $K_0$. The function~$f$ which maps $x$ to
$\tuple{x, x}$ is an example of a \emph{reduction} of $K$ to $K_0$.
\end{explain}

\begin{defn}
Let $A$ and $B$ be sets of natural numbers. A computable
function~$f\colon \Nat \to \Nat$ is a \emph{many-one reduction} of $A$
to~$B$ iff, for every natural number~$x$,
\[
x \in A \quad \text{if and only if} \quad f(x) \in B.
\]
If such a reduction $f$ exists, we say that $A$ is \emph{many-one
reducible} to~$B$, written $A \leq_m B$. If $A$ is many-one reducible
to $B$ and vice-versa, then $A$ and $B$ are said to be \emph{many-one
equivalent}, written $A \equiv_m B$.
\end{defn}

If the function $f$ in the definition above happens to be injective,
$A$~is said to be \emph{one-one reducible} to~$B$. Most of the
reductions described below meet this stronger requirement, but we will
not use this fact.

\begin{digress}
It is true, but by no means obvious, that one-one reducibility really
is a stronger requirement than many-one reducibility. In other words,
there are infinite sets $A$ and~$B$ such that $A$ is many-one
reducible to~$B$ but not one-one reducible to~$B$.
\end{digress}


% ==== OLP file: content/computability/computability-theory/rice-theorem.tex ====
% Part: computability
% Chapter: computability-theory
% Section: rice-theorem



\olfileid{cmp}{thy}{rce}
\olsection{Rice's Theorem}

If you think about it, you will see that the specifics of $\fn{Tot}$
do not play into the proof of \olref[tot]{prop:total}. We designed
$h(x,y)$ to act like the constant function $j(y) = 0$ exactly when $x$
is in $K$; but we could just as well have made it act like any other
partial computable function under those circumstances. This
observation lets us state a more general theorem, which says, roughly,
that no nontrivial property of computable functions is decidable.

Keep in mind that $\cfind{0}$, $\cfind{1}$, $\cfind{2}$,~\dots is our
standard enumeration of the partial computable functions.

\begin{thm}[Rice's Theorem]
  Let $C$ be any set of partial computable functions, and let $A =
  \Setabs{n}{\cfind{n} \in C}$. If $A$ is computable, then either $C$
  is empty or $C$ is the set of all the partial computable
  functions.
\end{thm}

An {\em index set} is a set $A$ with the property that if $n$ and $m$
are indices which ``compute'' the same function, then either both $n$
and $m$ are in $A$, or neither is. It is not hard to see that the
set~$A$ in the theorem has this property. Conversely, if $A$~is an
index set and $C$~is the set of functions computed by these indices,
then $A = \Setabs{n}{\cfind{n} \in C}$.

\begin{explain}
With this terminology, Rice's theorem is equivalent to saying that no
nontrivial index set is decidable. To understand what the theorem
says, it is helpful to emphasize the distinction between
\emph{programs} (say, in your favorite programming language) and the
functions they compute. There are certainly questions about programs
(indices), which are syntactic objects, that are computable: does this
program have more than 150 symbols? Does it have more than 22 lines?
Does it have a ``while'' statement? Does the string ``hello world''
ever appear as the argument to a ``print'' statement? Rice's theorem
says that no nontrivial question about the program's \emph{behavior}
is computable. This includes questions like these: does the program
halt on input $0$? Does it ever halt? Does it ever output an even
number?
\end{explain}

\begin{proof}[Proof of Rice's theorem]
Suppose $C$ is neither empty nor the set of all the partial
computable functions, and let $A$ be the set of indices of functions
in~$C$. We will show that if $A$ were computable, we could solve the
halting problem; so $A$~is not computable.

Without loss of generality, we can assume that the function $f$ which
is nowhere defined is not in $C$ (otherwise, switch $C$ and its
complement in the argument below). Let $g$ be any function in~$C$. The
idea is that if we could decide~$A$, we could tell the difference
between indices computing~$f$, and indices computing~$g$; and then we
could use that capability to solve the halting problem.

Here's how. Using the universal partial computable functions, we can
define a function
\[
h(x,y) \simeq
\begin{cases}
\text{undefined} & \text{if $\cfind{x}(x) \fundefined$} \\
g(y) & \text{otherwise.}
\end{cases}
\]
To compute $h$, first we try to compute $\cfind{x}(x)$; if that
computation halts, we go on to compute~$g(y)$; and if {\em that}
computation halts, we return the output. More formally, we can write
\[
h(x,y) \simeq \Proj{2}{0}(g(y),\fn{Un}(x,x)).
\]
where $\Proj{2}{0}(z_0, z_1) = z_0$ is the $2$-place projection
function returning the $0$-th argument, which is computable.

Then $h$ is a composition of partial computable functions, and the right
side is defined and equal to~$g(y)$ just when $\fn{Un}(x,x)$ and
$g(y)$ are both defined.

Notice that for a fixed~$x$, if $\cfind{x}(x)$ is undefined, then
$h(x,y)$ is undefined for every~$y$; and if $\cfind{x}(x)$ is defined,
then $h(x,y) \simeq g(y)$. So, for any fixed value of~$x$, either
$h(x,y)$ acts just like $f$ or it acts just like $g$, and deciding
whether or not $\cfind{x}(x)$ is defined amounts to deciding which of
these two cases holds. But this amounts to deciding whether or not
$h_x(y) \simeq h(x,y)$ is in~$C$ or not, and if $A$ were computable,
we could do just that.

More formally, since $h$~is partial computable, it is equal to the
function $\cfind{e}$ for some index~$e$. By the $s$-$m$-$n$ theorem
there is a primitive recursive function~$s$ such that for each~$x$,
$\cfind{s(e,x)}(y) = h_x(y)$. Now we have that for each $x$, if
$\cfind{x}(x) \fdefined$, then $\cfind{s(e,x)}$ is the same function
as~$g$, and so $s(e,x)$ is in $A$. On the other hand, if $\cfind{x}(x)
\uparrow$, then $\cfind{s(e,x)}$ is the same function as~$f$, and so
$s(e,x)$ is not in~$A$. In other words, we have that for every~$x$, $x
\in K$ if and only if $s(e,x) \in A$. If $A$ were computable,
$K$~would be also, which is a contradiction. So $A$ is not computable.
\end{proof}

Rice's theorem is very powerful. The following immediate corollary
shows some sample applications.
\begin{cor}
The following sets are undecidable.
\begin{enumerate}
\item $\Setabs{x}{\text{$17$ is in the range of $\cfind{x}$}}$
\item $\Setabs{x}{\text{$\cfind{x}$ is constant}}$
\item $\Setabs{x}{\text{$\cfind{x}$ is total}}$
\item $\Setabs{x}{\text{whenever $y < y'$, $\cfind{x}(y) \fdefined$, and
    if $\cfind{x}(y') \fdefined$, then $\cfind{x}(y) < \cfind{x}(y')$}}$
\end{enumerate}
\end{cor}

\begin{proof}
  These are all nontrivial index sets.
\end{proof}



% ==== OLP file: content/computability/computability-theory/russells-paradox.tex ====
% Part: computability
% Chapter: computability-theory
% Section: russells-paradox



\olfileid{cmp}{thy}{rus}
\olsection{Comparison with Russell's Paradox}

It is instructive to compare and contrast the arguments in
this section with Russell's paradox:
\begin{enumerate}
\item Russell's paradox: let $S = \Setabs{x}{x \notin x}$. Then $S
  \in S$ if and only if $S \notin S$, a contradiction.

  \emph{Conclusion:} There is no such set~$S$. Assuming the existence of a
  ``set of all sets'' is inconsistent with the other axioms of set
  theory.

\item A modification of Russell's paradox: let $F$ be the ``function''
  from the set of all functions to $\{ 0, 1 \}$, defined by
  \[
  F(f) =
  \begin{cases}
    1 & \text{if $f$ is in the domain of $f$, and $f(f) = 0$} \\
    0 & \text{otherwise}
  \end{cases}
  \]
  A similar argument shows that $F(F) = 0$ if and only if $F(F) = 1$,
  a contradiction.

  \emph{Conclusion:} $F$ is not a function. The ``set of all
  functions'' is too big to be the domain of a function.

\item The diagonalization argument: let $f_0$, $f_1$, \dots be the
  enumeration of the partial computable functions, and let $G \colon \Nat \to
  \{ 0, 1 \}$ be defined by
  \[
  G(x) =
  \begin{cases}
    1 & \text{if $f_x(x)\downarrow = 0$} \\
    0 & \text{otherwise}
  \end{cases}
  \]
  If $G$ is computable, then it is the function $f_k$ for some
  $k$. But then $G(k) = 1$ if and only if $G(k) = 0$, a contradiction.

  \emph{Conclusion:} $G$ is not computable. Note that according to the
  axioms of set theory, $G$ is still a function; there is no paradox
  here, just a clarification.
\end{enumerate}

That talk of partial functions, computable functions,
partial computable functions, and so on can be confusing. The set of
all partial functions from $\Nat$ to $\Nat$ is a big collection of
objects. Some of them are total, some of them are computable, some are
both total and computable, and some are neither. Keep in mind that
when we say ``function,'' by default, we mean a total function. Thus we
have:
\begin{enumerate}
\item computable functions
\item partial computable functions that are not total
\item functions that are not computable
\item partial functions that are neither total nor computable
\end{enumerate}
To sort this out, it might help to draw a big square representing all
the partial functions from $\Nat$ to $\Nat$, and then mark off two
overlapping regions, corresponding to the total functions and the
computable partial functions, respectively. It is a good exercise to
see if you can describe an object in each of the resulting regions in
the diagram.


% ==== OLP file: content/computability/computability-theory/s-m-n.tex ====
% Part: computability
% Chapter: computability-theory
% Section: s-m-n



\olfileid{cmp}{thy}{smn}
\olsection{The $s$-$m$-$n$ Theorem}

\begin{explain}
The next theorem is known as the ``$s$-$m$-$n$ theorem,'' for a reason
that will be clear in a moment. The hard part is understanding just
what the theorem says; once you understand the statement, it will seem
fairly obvious.
\end{explain}

\begin{thm}
\ollabel{thm:s-m-n}
  For each pair of natural numbers $n$ and~$m$, there is a primitive
  recursive function~$s^m_n$ such that for every sequence
  $e$, $a_0$, \dots, $a_{m-1}$, $y_0$ ,\dots, $y_{n-1}$, we have
  \[
  \cfind{s^m_n(e, a_0, \dots, a_{m-1})}[n](y_0, \dots, y_{n-1}) \simeq
  \cfind{e}[m+n](a_0, \dots, a_{m-1}, y_0, \dots, y_{n-1}).
\]
\end{thm}

\begin{explain}
It is helpful to think of $s^m_n$ as acting on \emph{programs}. That
is, $s^m_n$ takes a program~$e$ for an $(m+n)$-ary function, as well
as fixed inputs $a_0$, \dots, $a_{m-1}$; and it returns a program
$s^m_n(x, a_0, \dots, a_{m-1})$ for the $n$-ary function of the
remaining arguments. \iftag{TMs}{It you think of $x$ as the description of a
Turing machine, then $s^m_n(e, a_0, \dots, a_{m-1})$ is the Turing
machine that, on input $y_0$, \dots,~$y_{n-1}$, prepends
$a_0$, \dots,~$a_{m-1}$ to the input string, and runs~$e$. Each $s^m_n$
is then just a primitive recursive function that finds a code for the
appropriate Turing machine.}{}
\end{explain}



% ==== OLP file: content/computability/computability-theory/total.tex ====
% Part: computability
% Chapter: computability-theory
% Section: total



\olfileid{cmp}{thy}{tot}
\olsection{Totality is Undecidable}

Let us consider one more example of using the $s$-$m$-$n$ theorem to show
that something is noncomputable. Let $\fn{Tot}$ be the set of indices
of total computable functions, i.e.
\[
\fn{Tot} = \Setabs{x}{\text{for every $y$, $\cfind{x}(y)\fdefined$}}.
\]

\begin{prop}
\ollabel{prop:total}
$\fn{Tot}$ is not computable.
\end{prop}

\begin{proof}
To see that $\fn{Tot}$ is not computable, it suffices to show that $K$
is reducible to it. Let $h(x,y)$ be defined by
\[
h(x,y) \simeq
\begin{cases}
0 & \text{if $x \in K$} \\
\fundefined & \text{otherwise}
\end{cases}
\]
Note that $h(x,y)$ does not depend on $y$ at all. It should
not be hard to see that $h$~is partial computable: on input $x, y$, the
we compute~$h$ by first simulating the function~$\cfind{x}$ on input~$x$; if
this computation halts, $h(x,y)$ outputs $0$ and halts. So
$h(x,y)$ is just $\Zero(\umin{s}{T(x,x,s)})$, where $\Zero$ is the constant zero
function.

Using the $s$-$m$-$n$ theorem, there is a primitive recursive
function~$k(x)$ such that for every $x$ and~$y$,
\[
\cfind{k(x)}(y) =
\begin{cases}
0 & \text{if $x \in K$} \\
\fundefined & \text{otherwise}
\end{cases}
\]
So $\cfind{k(x)}$ is total if $x \in K$, and undefined otherwise. Thus,
$k$ is a reduction of $K$ to~$\fn{Tot}$.
\end{proof}

\begin{digress}
It turns out that $\fn{Tot}$ is not even computably enumerable---its
complexity lies further up on the ``arithmetical hierarchy.''  But we
will not worry about this strengthening here.
\end{digress}



% ==== OLP file: content/computability/computability-theory/universal-part-function.tex ====
% Part: computability
% Chapter: computability-theory
% Section: universal-part-function



\olfileid{cmp}{thy}{uni}
\olsection{The Universal Partial Computable Function}

\begin{thm}
\ollabel{thm:univ-comp}
There is a universal partial computable function~$\fn{Un}(e,x)$. In other
words, there is a function $\fn{Un}(e,x)$ such that:
\begin{enumerate}
\item $\fn{Un}(e,x)$ is partial computable.
\item If $f(x)$ is any partial computable function, then there is a
natural number $e$ such that $f(x) \simeq \fn{Un}(e,x)$ for every~$x$.
\end{enumerate}
\end{thm}

\begin{proof}
Let $\fn{Un}(e,x) \simeq U(\umin{s}{T(e,x,s)})$, where $U$ and $T$ are
as in Kleene's normal form theorem (\olref[nfm]{thm:normal-form}).
\end{proof}

\begin{explain}
This is just a precise way of saying that we have an effective
enumeration of the partial computable functions; the idea is that if
we write $f_e$ for the function defined by $f_e(x) = \fn{Un}(e,x)$,
then the sequence $f_0$, $f_1$, $f_2$, \dots includes all the partial
computable functions, with the property that $f_e(x)$ can be computed
``uniformly'' in $e$ and~$x$. For simplicity, we are using a binary
function that is universal for unary functions, but by coding
sequences of numbers we can easily generalize this to more arguments.
For example, note that if $f(x,y,z)$ is a $3$-place partial recursive
function, then the function $g(x) \simeq f((x)_0, (x)_1, (x)_2)$ is a
unary recursive function.
\end{explain}


